\documentclass[10pt,a4paper]{amsart}
\usepackage[margin=19mm]{geometry}
\usepackage{amsmath,amssymb,mathtools}
\usepackage[scr=rsfs,cal=euler]{mathalfa}
\usepackage{xfrac}
\usepackage[unicode,hidelinks,bookmarks=false]{hyperref}
\newcommand{\ie}{i.e.\ }
\newcommand{\cf}{cf.\ }
\newcommand{\resp}{respectively\ }
\newcommand{\Subsection}{\subsection}
\providecommand{\nobreakdash}{\nobreak\hbox{-}}
\setlength{\emergencystretch}{2em}
\multlinegap0pt
\usepackage{amsthm}
\newtheorem{theorem}{Theorem}
\newtheorem{claim}[theorem]{Claim}
\newtheorem{lemma}[theorem]{Lemma}
\DeclareMathOperator{\sing}{Sing}
\title[On moduli of foliated surfaces]{On moduli of foliated surfaces (Lemma 4.8)}
\author{Calum Spicer, Roberto Svaldi, Sebastian Velazquez}
\date{Source: arXiv:2511.13491v2 (CC BY 4.0)}
\begin{document}
\maketitle

\noindent\emph{Source and licence.} On moduli of foliated surfaces. Version arXiv:2511.13491v2 (CC BY 4.0), distributed under the Creative Commons Attribution 4.0 International licence, as recorded in the arXiv licence metadata for this exact version and shown on the abstract page https://arxiv.org/abs/2511.13491v2. Full rights notice: licenses/arxiv-2511.13491-CC-BY-4.0.txt.\par\smallskip

\noindent\textit{Editorial scope.} Counted unit: the Lemma printed as Lemma 4.8 of the article (source0.tex lines 3081--3101, source label \texttt{lem\_sing\_adjunction\_families}), delivered together with the article's own complete original proof of it (source0.tex lines 3102--3203, one \texttt{proof} environment of 102 lines). No mathematics is written, repaired or re-derived: statement and proof are the article's own text line for line. Required context retained uncounted: none. Every definition, display and label of those lines is retained unchanged. Omitted from this package and not redistributed: 1--3080, 3204--4967. Nothing inside a retained excerpt is omitted or altered: every retained body is delivered exactly as the article's lines are written, so no omission receipt of any kind accompanies this package. Citations: the retained excerpts carry no citation command at all, so no bibliography entry of the article is transcribed and no reference is redistributed. Reference adaptation: none was needed and none was made; every reference inside the delivered unit resolves inside it, so no unresolved reference or citation is produced. Printed slips kept literal: no printed word, symbol, hypothesis or number of the counted statement or of its proof was altered. Layout and class: the article's own class is \texttt{amsart} with packages including inputenc, fancyhdr, geometry, hyperref, subcaption, float, graphicx, graphics, color, tikz-cd, tikz, listings, cite, enumerate; the wrapper is the lane's pinned \texttt{amsart} editorial wrapper at 10pt on A4, which restates the theorem-like environments the retained text needs and adds the article's own macro definitions that the retained text needs (\texttt{\textbackslash sing}), with extra wrapper packages (none). The delivered numbering is the unit's own. Assets: no retained span contains a figure environment, a graphics-inclusion command, a PDF-page inclusion, a table or a TikZ picture; no image file is copied into this package, the package declares no asset dependency and no image file is redistributed with it. Licence: version arXiv:2511.13491v2 of this article is distributed under the Creative Commons Attribution 4.0 International licence, as recorded in the arXiv licence metadata for this exact version and shown on the abstract page https://arxiv.org/abs/2511.13491v2. A full rights notice is delivered at licenses/arxiv-2511.13491-CC-BY-4.0.txt. Characters: every retained excerpt is pure ASCII, so the delivered engine, XeLaTeX with its default Latin Modern text font, reports no missing character.
\begin{lemma}
	\label{lem_sing_adjunction_families}

Consider the following set up.
\begin{enumerate}
    \item Let $X$ be a smooth variety with $\dim X = 3$. 
    \item Let $f\colon X \to T$ be a surjective morphism to a smooth curve with simple normal crossings fibres.
    \item Let $\mathcal F$ be a rank one integrable distribution on $X$ such that
    \begin{enumerate} 
    \item $T_{\mathcal F} \subset T_{X/T}$; 
    \item $\sing \mathcal F$ does not 
contain any components of fibres of $f$; and
\item  $\mathcal F$ has log canonical singularities.
\end{enumerate}
\end{enumerate}

	Let $0 \in T$ be a closed point
	and let $S_0 \subset X_0 = f^{-1}(0)$ be an irreducible component. 
	Let $(\mathcal F_0, D_0)$ be the restricted foliated pair on $S_0$. 
Then for any component $C$ of $D_0$ we have $m_CD_0 \leq \varepsilon(\mathcal F_0, C)$.
\end{lemma}

\begin{proof}
	Let $C$ be an irreducible component of $D_0$.  The question is local about a general point of  $C$, so we may freely choose a system of coordinates $(t, x, y)$ on $X$ so that
	 $C = \{x = t = 0\}$, $S_0 = \{t = 0\}$ and 
     \begin{itemize}
         \item in the case that $X_0$ is normal, $f(t, x, y) = t^s$; and 
         \item  in the case that $X_0$ is not normal in a neighbourhood of $C$, $f(t, x, y) = t^sx^r$, so $(X_0)_{{\rm red}} = \{xt = 0\}$.
     \end{itemize}
 
	We may assume that $\mathcal F$ is defined by a vector field  $\partial = a\partial_x+b\partial_t+c\partial_y$
    where $a, b, c \in \mathcal O_X$.  Since $S_0$ is invariant by $\partial$ let us note that $t|b$.



Let us suppose for sake of contradiction that $m_CD_0 = k >\varepsilon(\mathcal F_0, C)$.
We remark that any exceptional divisor centred over $C$ is mapped to $0 \in T$, and is therefore invariant by the transform of $\mathcal F$.  Thus to derive a contradiction it suffices to find an exceptional divisor $E$ dominating $C$ with discrepancy $a(E, \mathcal F) \leq -1 < 0 = \varepsilon(E)$.




Note that $m_CD_0$ is precisely the order of vanishing of the restricted vector field $\partial|_{\{t = 0\}}$ along $C$, in particular, $C \subset \sing \partial$.  Thus,
$m_CD_0 = k$
implies that $a, b, c \in (x^k, t)$.
So, let us write $a = A_1x^k+A_2t, b = B_1x^k+B_2t$
and $c = C_1x^k+C_2t$ where $A_i, B_i, C_i \in \mathcal O_X$.
As noted earlier, $t|b$, so in fact we may assume that $B_1 = 0$.


\begin{claim}
\label{claim_order2}
$A_1x^k \in (x, t)^2$.
\end{claim}
\begin{proof}[Proof of Claim]
The restriction of $\partial$ to $S_0$ is 
$x^k\alpha_1\partial_x+x^k\gamma_1\partial_y$ where $\alpha_1$ (resp. $\gamma_1$) is the restriction of $A_1$ (resp. $C_1$) to $S_0$. This vector field has a codimension $=1$ component of its zero locus.  After dividing by an appropriate function we may find a generator of $\mathcal F_0$.  Indeed, we may write $x^k\alpha_1 = x^kq\tilde{\alpha}_1$ and 
$x^k\gamma_1 = x^kq\tilde{\gamma}_1$
where $\tilde{\alpha}_1$ and $\tilde{\gamma}_1$ have no common factors.  In particular, $\tilde{\alpha}_1\partial_x+\tilde{\gamma}_1\partial_y$
is a vector field on $S_0$ without codimension one zeros which defines $\mathcal F_0$.

We argue in cases based on $\varepsilon(\mathcal F_0, C)$.

If $\varepsilon(\mathcal F_0, C) = 1$, then $k \ge 2$, and so $A_1x^k \in (x, t)^2$.  

If $\varepsilon(\mathcal F_0, C) = 0$ then $k \ge 1$.   Since $\{x = 0\}$ is $\mathcal F_0$-invariant, it is invariant by the vector field 
$\tilde{\alpha}_1\partial_x+\tilde{\gamma}_1\partial_y$.  It follows that $x|\tilde{\alpha}_1$, hence $x|\alpha_1$, which implies that $A_1 = xA'_1+tA''_1$ where $A'_1, A''_1 \in \mathcal O_X$, and so $A_1x^k \in (x, t)^2$.
\end{proof}

We next make the following observation.

\begin{claim}
\label{claim_easy_discrep_calc}
    If
$a, b\in (x, t)^2$, 
then the blow up in $C$ will extract a divisor of discrepancy $\leq -1$.
\end{claim}
\begin{proof}[Proof of Claim]
Consider the coordinate chart for the blow up $p\colon X' \rightarrow X$ in $C$ 
given by $(u, v, y) \mapsto (uv, u, y) = (t, x, y)$.
In these coordinates we can verify that 
$p^*\partial_x = \partial_u - \frac{v}{u}\partial_v, p^*\partial_t = \frac{1}{u}\partial_v$ and $p^*\partial_y = \partial_y$ and so 
\begin{align*}p^*((A_1x^k+A_2t)\partial_x+B_2t\partial_t+(C_1x^k+C_2t)\partial_y) =& \\
(\tilde A_1u^k+\tilde A_2uv)\partial_u +(\tilde B_2v-\tilde A_1vu^{k-1}-\tilde A_2v^2)&\partial_v+(\tilde C_1u^k+\tilde C_2uv)\partial_y,
\end{align*}
where we denote by $\tilde{A}_i$ (resp. $\tilde{B}_i$, resp. $\tilde{C}_i$) the pullback of $A_i$ (resp. $B_i$, resp. $C_i$)
by $p$).
Since $a = A_1x^k+A_2t, b = B_2t \in (x, t)^2$ and $A_1x^k \in (x, t)^2$ (by Claim \ref{claim_order2}), it follows that
$A_2, B_2 \in (x, t)$, hence $u|\tilde A_2$ and 
$u|\tilde B_2$. Since $A_1x^k \in (x, t)^2$ we see that $u^2|\tilde A_1u^k$ and so $u| (\tilde B_2v-\tilde A_1vu^{k-1}-\tilde A_2v^2)$.  It follows that $p^*\partial$ vanishes along $\{u = 0\}$, i.e., $p$ extracts a divisor of discrepancy $\le-1$.
\end{proof}

We argue in cases based on whether $X_0$ is normal at the generic point of $C$ or not.

{\bf Case I: }{\it  Suppose that $X_0$ is normal at the generic point of $C$.}

In this case we must have $\partial(t) = 0$ 
and so $b = 0$.
Consider the coordinate chart for the blow up $p\colon X' \rightarrow X$ in $C$ 
given by $(u, v, y) \mapsto (uv, u, y) = (t, x, y)$.
In these coordinates we can verify that 
\begin{align*}p^*((A_1x^k+A_2t)\partial_x+(C_1x^k+C_2t)\partial_y)) =& \\
(\tilde A_1u^k+\tilde A_2uv)\partial_u - &(\tilde A_1vu^{k-1}+\tilde A_2v^2)\partial_v+(\tilde C_1u^k+\tilde C_2uv)\partial_y
\end{align*}
in particular this blow up extracts an invariant divisor 
of discrepancy $\leq 0$.
By Claim \ref{claim_order2} $A_1x^k \in (x, t)^2$ and so 
we see that 
\begin{align*}(\tilde A_1u^k+\tilde A_2uv), (\tilde A_1vu^{k-1}+\tilde A_2v^2)
\in (u, v)^2.\end{align*}  By Claim \ref{claim_easy_discrep_calc} a blow up centred in $\{u = v = 0\}$
will extract a divisor of discrepancy $\leq -1$, a contradiction.


\medskip 

{\bf Case II: }{\it  Suppose that $X_0$ is not normal at the generic point of $C$.}

In this case $X_0 = \{x^rt^s = 0\}$ where $r, s \in \mathbb N$ and so 
\begin{align*}0 = \partial(x^rt^s) = rx^{r-1}t^sa+sx^rt^{s-1}b.\end{align*}
In particular, we see that $x|a$ and $t|b$.
Since $x|a$ we see that $x|A_2$ and since $A_1x^k \in (x, t)^2$ by Claim \ref{claim_order2} we see that $a\in (x, t)^2$.
Since $rta+sxb = 0$ we then deduce that $b \in (x, t)^2$ and so 
Claim \ref{claim_easy_discrep_calc} shows that
the blow up in $\{x = t = 0\}$ extracts a divisor of discrepancy $\le-1$, a contradiction.
\end{proof}

\end{document}
